\chapter{2005年山东卷（文）}
\section{单选题}
\begin{problem}
\(\{a_n\}\) 是首项 \(a_1 = 1\)，公差 \(d = 3\) 的等差数列. 如果 \(a_n = 2005\)，则序号 \(n\) 等于（\(\quad\)）

\choices
    {\(667\)}
    {\(668\)}
    {\(669\)}
    {\(670\)}
\end{problem}

\begin{answer}
C
\end{answer}

\begin{solution}
由等差数列通项公式，\(a_n=a_1+(n-1)d=1+3(n-1)=3n-2\). 由 \(a_n=2005\) 得 \(3n-2=2005\)，所以 \(n=669\)，故选 C.
\end{solution}

\begin{problem}
下列大小关系正确的是（\(\quad\)）

\choices
    {\(0.4^{3} < 3^{0.4} < \log_{4}0.3\)}
    {\(0.4^{3} < \log_{4}0.3 < 3^{0.4}\)}
    {\(\log_{4}0.3 < 0.4^3 < 3^{0.4}\)}
    {\(\log_{4}0.3 < 3^{0.4} < 0.4^3\)}
\end{problem}

\begin{answer}
C
\end{answer}

\begin{solution}
因为 \(0<0.4<1\)，所以 \(0<0.4^3<1\)；又因为 \(3>1\)，所以 \(3^{0.4}>1\). 由于对数的底数 \(4>1\)，而 \(0<0.3<1\)，故 \(\log_4 0.3<0\). 因此 \(\log_4 0.3<0.4^3<3^{0.4}\)，故选 C.
\end{solution}

\begin{problem}
函数 \( y = \frac{1 - x}{x} \)（\(x \neq 0\)）的反函数图象大致是（\(\quad\)）

\choices
    {\choicebitmap[width=0.15\paperwidth]{img/2005/shandong_liberal/q3_optA_nolabel.png}}
    {\choicebitmap[width=0.15\paperwidth]{img/2005/shandong_liberal/q3_optB_nolabel.png}}
    {\choicebitmap[width=0.15\paperwidth]{img/2005/shandong_liberal/q3_optC_nolabel.png}}
    {\choicebitmap[width=0.15\paperwidth]{img/2005/shandong_liberal/q3_optD_nolabel.png}}
\end{problem}

\begin{answer}
B
\end{answer}

\begin{solution}
由 \(y=\dfrac{1-x}{x}\) 得 \(x(y+1)=1\)，所以反函数为 \(y=\dfrac{1}{x+1}\)，其定义域为 \(x\neq -1\). 该双曲线的中心为 \((-1,0)\)，且当 \(x>-1\) 时 \(y>0\)，当 \(x<-1\) 时 \(y<0\). 因此图象与选项 B 相符，故选 B.
\end{solution}

\begin{problem}
已知函数 \( y = \sin \left( x - \frac{\pi}{12} \right) \cos \left( x - \frac{\pi}{12} \right) \)，则下列判断正确的是（\(\quad\)）

\choices
    {此函数的最小周期为 \(2\pi\)，其图象的一个对称中心是 \(\left(\frac{\pi}{12}, 0\right)\)}
    {此函数的最小周期为 \(\pi\)，其图象的一个对称中心是 \(\left(\frac{\pi}{12}, 0\right)\)}
    {此函数的最小周期为 \(2\pi\)，其图象的一个对称中心是 \(\left(\frac{\pi}{6}, 0\right)\)}
    {此函数的最小周期为 \(\pi\)，其图象的一个对称中心是 \(\left(\frac{\pi}{6}, 0\right)\)}
\end{problem}

\begin{answer}
B
\end{answer}

\begin{solution}
令 \(u=x-\dfrac{\pi}{12}\)，则 \(y=\sin u\cos u=\dfrac12\sin 2u=\dfrac12\sin\left(2x-\dfrac{\pi}{6}\right)\). 因此最小正周期为 \(\pi\). 当 \(x=\dfrac{\pi}{12}\) 时，\(u=0\)，函数值为 \(0\)，且函数关于该点中心对称，所以一个对称中心为 \(\left(\dfrac{\pi}{12},0\right)\)，故选 B.
\end{solution}

\begin{problem}
下列函数既是奇函数，又在区间 \([-1, 1]\) 上单调递减的是（\(\quad\)）

\choices
    {\(f(x) = \sin x\)}
    {\(f(x) = -|x + 1|\)}
    {\(f(x) = \frac{1}{2} (a^x + a^{-x})\)}
    {\(f(x) = \ln \frac{2 - x}{2 + x}\)}
\end{problem}

\begin{answer}
D
\end{answer}

\begin{solution}
选项 A 中，\(f'(x)=\cos x>0\)（\(-1\leqslant x\leqslant1<\dfrac\pi2\)），所以函数在 \([-1,1]\) 上单调递增，不合题意. 选项 B 中，\(f(0)=-1\neq0\)，不可能是奇函数. 选项 C 满足 \(f(-x)=f(x)\)，是偶函数而不是奇函数.

选项 D 的定义域包含 \([-1,1]\)，且
\(f(-x)=\ln\dfrac{2+x}{2-x}=-\ln\dfrac{2-x}{2+x}=-f(x)\)，所以它是奇函数. 又
\(f'(x)=-\dfrac{1}{2-x}-\dfrac{1}{2+x}=-\dfrac{4}{4-x^2}<0\)（\(-1\leqslant x\leqslant1\)），故它在 \([-1,1]\) 上单调递减. 因此选 D.
\end{solution}

\begin{problem}
如果 \(\left(3x - \frac{1}{\sqrt[3]{x^2}}\right)^n\) 的展开式中各项系数之和为 \(128\)，则展开式中 \(\frac{1}{x^3}\) 的系数是（\(\quad\)）

\choices
    {\(7\)}
    {\(-7\)}
    {\(21\)}
    {\(-21\)}
\end{problem}

\begin{answer}
C
\end{answer}

\begin{solution}
令展开式中的 \(x=1\)，各项系数之和为 \(\left(3-1\right)^n=2^n\). 所以 \(2^n=128=2^7\)，得 \(n=7\). 展开式的通项中，若从第二项取 \(k\) 次，则该项为
\(\mathrm C_7^k3^{7-k}(-1)^kx^{7-k-\frac{2k}{3}}\). 要含有 \(x^{-3}\)，应满足 \(7-k-\dfrac{2k}{3}=-3\)，即 \(k=6\). 所求系数为 \(\mathrm C_7^6 3=21\)，故选 C.
\end{solution}

\begin{problem}
函数 \( f(x) = \begin{cases}
\sin (\pi x^2), & -1 < x < 0,\\
\e^{x - 1}, & x \geqslant 0.
\end{cases} \) 若 \( f(1) + f(a) = 2 \)，则 \( a \) 的所有可能值为（\(\quad\)）

\choices
    {\(1\)}
    {\( -\frac{\sqrt{2}}{2} \)}
    {\(1\)，\(-\frac{\sqrt{2}}{2}\)}
    {\(1\)，\(\frac{\sqrt{2}}{2}\)}
\end{problem}

\begin{answer}
C
\end{answer}

\begin{solution}
由 \(f(1)=\e^0=1\)，得 \(f(a)=1\). 当 \(-1<a<0\) 时，\(0<a^2<1\)，所以 \(0<\pi a^2<\pi\)，方程 \(\sin(\pi a^2)=1\) 只有解 \(\pi a^2=\dfrac\pi2\)，从而 \(a=-\dfrac{\sqrt2}{2}\). 当 \(a\geqslant0\) 时，\(\e^{a-1}=1\)，所以 \(a=1\). 因此所有可能值为 \(1\)，\(-\dfrac{\sqrt2}{2}\)，故选 C.
\end{solution}

\begin{problem}
已知向量 \(\bs{a}\)，\(\bs{b}\)，且 \(\overrightarrow{AB} = \bs{a} + 2\bs{b}\)，\(\overrightarrow{BC} = -5\bs{a} + 6\bs{b}\)，\(\overrightarrow{CD} = 7\bs{a} - 2\bs{b}\)，则一定共线的三点是（\(\quad\)）

\choices
    {A，B，D}
    {A，B，C}
    {B，C，D}
    {A，C，D}
\end{problem}

\begin{answer}
A
\end{answer}

\begin{solution}
由向量的加法，
\(\overrightarrow{AD}=\overrightarrow{AB}+\overrightarrow{BC}+\overrightarrow{CD}=3\bs{a}+6\bs{b}=3(\bs{a}+2\bs{b})=3\overrightarrow{AB}\). 因而 \(\overrightarrow{AD}\) 与 \(\overrightarrow{AB}\) 平行，点 \(A\)，\(B\)，\(D\) 共线，故选 A.
\end{solution}

\begin{problem}
设地球的半径为 \(R\)，若甲地位于北纬 \( 45^{\circ} \) 东经 \( 120^{\circ} \)，乙地位于南纬 \( 75^{\circ} \) 东经 \( 120^{\circ} \)，则甲，乙两地的球面距离为（\(\quad\)）

\choices
    {\(\sqrt{3} R\)}
    {\(\frac{5}{6} R\)}
    {\(\frac{5\pi}{6} R\)}
    {\(\frac{2\pi}{3} R\)}
\end{problem}

\begin{answer}
D
\end{answer}

\begin{solution}
甲、乙两地经度相同，所以它们在同一条经线上. 甲地纬度为 \(45^{\circ}\)，乙地纬度为 \(-75^{\circ}\)，两地对应的地心角为 \(45^{\circ}-(-75^{\circ})=120^{\circ}=\dfrac{2\pi}{3}\). 因此球面距离为 \(R\cdot\dfrac{2\pi}{3}=\dfrac{2\pi}{3}R\)，故选 D.
\end{solution}

\begin{problem}
\(10\) 张奖券中只有 \(3\) 张有奖，\(5\) 个人购买，每人 \(1\) 张. 至少有 \(1\) 人中奖的概率是（\(\quad\)）

\choices
    {\(\frac{3}{10}\)}
    {\(\frac{1}{12}\)}
    {\(\frac{1}{2}\)}
    {\(\frac{11}{12}\)}
\end{problem}

\begin{answer}
D
\end{answer}

\begin{solution}
用对立事件计算. 五个人都没有中奖的概率为 \(\dfrac{\mathrm C_7^5}{\mathrm C_{10}^5}=\dfrac{21}{252}=\dfrac1{12}\). 所以至少有一人中奖的概率为 \(1-\dfrac1{12}=\dfrac{11}{12}\)，故选 D.
\end{solution}

\begin{problem}
设集合 \(A\)，\(B\) 是全集 \(U\) 的两个子集，则 \(A \subsetneq B\) 是 \((\complement_U A) \cup B = U\) 的（\(\quad\)）

\choices
    {充分不必要条件}
    {必要不充分条件}
    {充要条件}
    {既不充分也不必要条件}
\end{problem}

\begin{answer}
A
\end{answer}

\begin{solution}
由集合运算，\((\complement_U A)\cup B=U\) 等价于不存在属于 \(A\) 但不属于 \(B\) 的元素，即等价于 \(A\subseteq B\). 因此 \(A\subsetneq B\) 能推出该等式，但该等式成立时也可能有 \(A=B\)，所以 \(A\subsetneq B\) 是充分不必要条件，故选 A.
\end{solution}

\begin{problem}
设直线 \(l: 2x + y + 2 = 0\) 关于原点对称的直线为 \(l'\)，若 \(l'\) 与椭圆 \(x^2 + \frac{y^2}{4} = 1\) 的交点为 \(A\)，\(B\)，点 \(P\) 为椭圆上的动点，则使 \(\triangle PAB\) 的面积为 \(\frac{1}{2}\) 的点 \(P\) 的个数为（\(\quad\)）

\choices
    {\(1\)}
    {\(2\)}
    {\(3\)}
    {\(4\)}
\end{problem}

\begin{answer}
B
\end{answer}

\begin{solution}
关于原点对称后，\(l'\) 的方程为 \(2x+y-2=0\). 将椭圆参数化为 \(x=\cos t\)，\(y=2\sin t\)，代入 \(l'\) 得 \(\cos t+\sin t=1\)，所以两个交点为 \((1,0)\) 和 \((0,2)\)，从而 \(|AB|=\sqrt5\). 若点 \(P\) 到直线 \(AB\) 的距离为 \(h\)，则 \(\triangle PAB\) 的面积为 \(\dfrac12\sqrt5h\). 由面积为 \(\dfrac12\) 得 \(h=\dfrac1{\sqrt5}\)，于是
\(\left|2x+y-2\right|=1\)，即点 \(P\) 在直线 \(2x+y=1\) 或 \(2x+y=3\) 上.

对于直线 \(2x+y=c\)，代入椭圆得
\(2x^2-cx+\dfrac{c^2}{4}-1=0\)，其判别式为 \(8-c^2\). 当 \(c=1\) 时判别式为 \(7>0\)，有两个交点；当 \(c=3\) 时判别式为 \(-1<0\)，没有交点. 所以满足条件的点 \(P\) 有 \(2\) 个，故选 B.
\end{solution}

\section{填空题}
\begin{problem}
某学校共有教师 \(490\) 人，其中不到 \(40\) 岁的有 \(350\) 人，\(40\) 岁及以上的有 \(140\) 人，为了解普通话在该校教师中的推广普及情况，用分层抽样的方法，从全体教师中抽取一个容量为 \(70\) 人的样本进行普通话水平测试，其中在不到 \(40\) 岁的教师中应抽取的人数是\fillinblank{}.
\end{problem}

\begin{answer}
50
\end{answer}

\begin{solution}
分层抽样时各层的抽样比相同. 不到 \(40\) 岁教师应抽取的人数为 \(70\cdot\dfrac{350}{490}=70\cdot\dfrac57=50\)，所以填 \(50\).
\end{solution}

\begin{problem}
设双曲线 \( \frac{x^{2}}{a^{2}} - \frac{y^{2}}{b^{2}} = 1 \)（\(a > 0\)，\(b > 0\)）的右焦点为 \(F\)，右准线 \(l\) 与两条渐近线交于 \(P\)，\(Q\) 两点，如果 \( \triangle PQF \) 是直角三角形，则双曲线的离心率 \( e = \)\fillinblank{}.
\end{problem}

\begin{answer}
\(\sqrt2\)
\end{answer}

\begin{solution}
设双曲线的焦半距为 \(c\)，则 \(c^2=a^2+b^2\)，右焦点为 \(F=(c,0)\)，右准线为 \(x=\dfrac{a^2}{c}\). 令 \(d=\dfrac{a^2}{c}\)，则渐近线与准线的交点为 \(P=\left(d,\dfrac{bd}{a}\right)\)，\(Q=\left(d,-\dfrac{bd}{a}\right)\). 因为 \(PQ\) 是竖直线，在 \(P\) 处两边的方向向量可取 \(\overrightarrow{PQ}=\left(0,-\dfrac{2ab}{c}\right)\)、\(\overrightarrow{PF}=\left(c-d,-\dfrac{ab}{c}\right)\)，有 \(\overrightarrow{PQ}\cdot\overrightarrow{PF}=\dfrac{2a^2b^2}{c^2}\ne0\). 在 \(Q\) 处，\(\overrightarrow{QP}=\left(0,\dfrac{2ab}{c}\right)\)、\(\overrightarrow{QF}=\left(c-d,\dfrac{ab}{c}\right)\)，所以 \(\overrightarrow{QP}\cdot\overrightarrow{QF}=\dfrac{2a^2b^2}{c^2}\ne0\). 因而直角只能在 \(F\) 处.

在 \(F\) 处，\(\overrightarrow{FP}=\left(d-c,\dfrac{ab}{c}\right)\)，\(\overrightarrow{FQ}=\left(d-c,-\dfrac{ab}{c}\right)\). 直角条件为
\(\overrightarrow{FP}\cdot\overrightarrow{FQ}=\left(c-d\right)^2-\left(\dfrac{ab}{c}\right)^2=0\)，即
\(\left(c-d\right)^2=\left(\dfrac{bd}{a}\right)^2\).

由于 \(c-d=\dfrac{c^2-a^2}{c}=\dfrac{b^2}{c}>0\)，而 \(\dfrac{bd}{a}=\dfrac{ab}{c}>0\)，取正平方根得 \(b^2=ab\). 由 \(a\)，\(b>0\) 得 \(a=b\)，从而 \(c=\sqrt2a\)，所以 \(e=\dfrac ca=\sqrt2\).
\end{solution}

\begin{problem}
设 \(x\)，\(y\) 满足约束条件\(\begin{cases}
x + y \leqslant 5,\\
3x + 2y \leqslant 12,\\
0 \leqslant x \leqslant 3,\\
0 \leqslant y \leqslant 4
\end{cases}\) 则使得目标函数 \(z = 6x + 5y\) 的最大的点 \((x,y)\) 是\fillinblank{}.
\end{problem}

\begin{answer}
\((2,3)\)
\end{answer}

\begin{solution}
由约束条件，
\(z=6x+5y=3(x+y)+(3x+2y)\leqslant3\cdot5+12=27\).

当且仅当 \(x+y=5\) 且 \(3x+2y=12\) 时取等号. 解方程组
\(\begin{cases}
x+y=5,\\
3x+2y=12
\end{cases}\) 得 \(x=2\)，\(y=3\)，且该点满足 \(0\leqslant x\leqslant3\)、\(0\leqslant y\leqslant4\). 因此目标函数取得最大值的点为 \((2,3)\)，所以填 \((2,3)\).
\end{solution}

\begin{problem}
已知 \(m\)，\(n\) 是不同的直线，\(\alpha\)，\(\beta\) 是不重合的平面，给出下列命题：

\begin{circlelist}
\item 若 \(m \parallel \alpha\)，则 \(m\) 平行于平面 \(\alpha\) 内的任一条直线；
\item 若 \(\alpha \parallel \beta\)，\(m \subset \alpha\)，\(n \subset \beta\)，则 \(m \parallel n\)；
\item 若 \(m \perp \alpha\)，\(n \perp \beta\)，\(m \parallel n\)，则 \(\alpha \parallel \beta\)；
\item 若 \(\alpha \parallel \beta\)，\(m \subset \alpha\)，则 \(m \parallel \beta\).
\end{circlelist}

上面的命题中，真命题的序号是（写出所有真命题的序号）\fillinblank{}.
\end{problem}

\begin{answer}
\(3\)，\(4\)
\end{answer}

\begin{solution}
命题 1 不成立，因为直线 \(m\) 虽与平面 \(\alpha\) 平行，但不一定与平面内的每一条直线平行. 命题 2 不成立，因为两个平行平面内任取的两条直线方向不一定相同. 命题 3 中，两条互相平行且分别垂直于两个平面的直线具有同一法向方向，因此两个不重合平面平行，命题 3 成立. 命题 4 中，平面 \(\alpha\) 内的直线不与平面 \(\beta\) 相交，且其方向平行于平面 \(\beta\)，所以命题 4 成立. 真命题的序号为 \(3\)，\(4\).
\end{solution}

\section{解答题}
\begin{problem}
已知向量 \(\bs{m} = (\cos \theta, \sin \theta)\) 和 \(\bs{n} = (\sqrt{2} - \sin \theta, \cos \theta)\)，\(\theta \in (\pi, 2\pi)\)，且 \(|\bs{m} + \bs{n}| = \frac{8\sqrt{2}}{5}\)，求 \(\cos \left(\frac{\theta}{2} + \frac{\pi}{8}\right)\) 的值.
\end{problem}

\begin{solution}
由题意，
\(\bs m+\bs n=(\cos\theta-\sin\theta+\sqrt2,\sin\theta+\cos\theta)\). 因此
\(\left|\bs m+\bs n\right|^2=4+2\sqrt2(\cos\theta-\sin\theta)\). 代入已知条件得
\(4+2\sqrt2(\cos\theta-\sin\theta)=\dfrac{128}{25}\)，从而 \(\cos\theta-\sin\theta=\dfrac{7\sqrt2}{25}\)，即
\(\cos\left(\theta+\dfrac\pi4\right)=\dfrac7{25}\).

因为 \(\theta\in(\pi,2\pi)\)，所以 \(\theta+\dfrac\pi4\in\left(\dfrac{5\pi}{4},\dfrac{9\pi}{4}\right)\). 令 \(\gamma=\arccos\dfrac7{25}\)，则在此区间内只能有 \(\theta+\dfrac\pi4=2\pi-\gamma\). 于是
\(\dfrac\theta2+\dfrac\pi8=\pi-\dfrac\gamma2\). 又 \(\sin\gamma=\dfrac{24}{25}\)，故
\(\cos\dfrac\gamma2=\sqrt{\dfrac{1+\cos\gamma}{2}}=\dfrac45\). 因而所求值为 \(-\dfrac45\).
\end{solution}

\begin{problem}
袋中装有黑球和白球共 \(7\) 个，从中任取 \(2\) 个球都是白球的概率为 \(\frac{1}{7}\)，现有甲，乙两人从袋中轮流摸取 \(1\) 球，甲先取，乙后取，然后甲再取…，取后不放回，直到两人中有一人取到白球时既终止，每个球在每一次被取出的机会是等可能的.

\begin{enumerate}
\item 求袋中原有的白球的个数；
\item 求取球 \(2\) 次终止的概率；
\item 求甲取到白球的概率.
\end{enumerate}
\end{problem}

\begin{solution}
\begin{enumerate}
\item 设白球有 \(w\) 个，则 \(\dfrac{\mathrm C_w^2}{\mathrm C_7^2}=\dfrac17\). 所以 \(w(w-1)=6\)，由 \(0\leqslant w\leqslant7\) 得 \(w=3\)，黑球有 \(4\) 个.
\item 取球两次终止必须先取到黑球、再取到白球，所以概率为 \(\dfrac47\cdot\dfrac36=\dfrac27\).
\item 甲在第 \(1\)，\(3\)，\(5\)，\(7\) 次取球. 白球在第 1 次被甲取到的概率为 \(\dfrac37\)；在第 3 次被甲取到的概率为 \(\dfrac47\cdot\dfrac36\cdot\dfrac35=\dfrac6{35}\)；在第 5 次被甲取到的概率为 \(\dfrac47\cdot\dfrac36\cdot\dfrac25\cdot\dfrac14\cdot1=\dfrac1{35}\). 第 7 次不可能才首次取到白球，因为只有 \(4\) 个黑球. 因此甲取到白球的概率为 \(\dfrac37+\dfrac6{35}+\dfrac1{35}=\dfrac{22}{35}\).
\end{enumerate}
\end{solution}

\begin{problem}
已知 \(x = 1\) 是函数 \(f(x) = mx^3 - 3(m + 1)x^2 + nx + 1\) 的一个极值点，其中 \(m\)，\(n \in \R\)，\(m < 0\).

\begin{enumerate}
\item 求 \(m\) 与 \(n\) 的关系式；
\item 求 \( f(x) \) 的单调区间.
\end{enumerate}
\end{problem}

\begin{solution}
\begin{enumerate}
\item 因为 \(x=1\) 是可导函数的极值点，所以 \(f'(1)=0\). 而 \(f'(x)=3mx^2-6(m+1)x+n\)，故 \(3m-6(m+1)+n=0\)，即 \(n=3m+6\).
\item 代入上式，得
\(f'(x)=3m\left(x-1\right)\left(x-\dfrac{m+2}{m}\right)\). 记 \(r=\dfrac{m+2}{m}=1+\dfrac2m\)，由于 \(m<0\)，有 \(r<1\). 又因 \(3m<0\)，所以 \(f'(x)<0\) 在 \((-\infty,r)\) 和 \((1,+\infty)\) 上成立，\(f'(x)>0\) 在 \((r,1)\) 上成立. 因此 \(f(x)\) 在 \(\left(-\infty,\dfrac{m+2}{m}\right)\) 和 \((1,+\infty)\) 上单调递减，在 \(\left(\dfrac{m+2}{m},1\right)\) 上单调递增.
\end{enumerate}
\end{solution}

\begin{problem}
如图，已知长方体 \(ABCD - A_{1}B_{1}C_{1}D_{1}\)，\(AB = 2\)，\(AA_{1} = 1\)，直线 \(BD\) 与平面 \(AA_{1}B_{1}B\) 所成的角为 \(30^{\circ}\)，\(AE\) 垂直 \(BD\) 于 \(E\)，\(F\) 为 \(A_{1}B_{1}\) 的中点.

\begin{enumerate}
\item 求异面直线 \(AE\) 与 \(BF\) 所成的角；
\item 求平面 \(BDF\) 与平面 \(AA_{1}B\) 所成的二面角（锐角）的大小；
\item 求点 \(A\) 到平面 \(BDF\) 的距离.
\end{enumerate}

\bitmapfigure[max width=0.46\linewidth]{img/2005/shandong_liberal/q22_fig1.png}
\end{problem}

\begin{solution}
建立直角坐标系，令 \(B=(0,0,0)\)，\(C=(w,0,0)\)，\(A=(0,2,0)\)，\(D=(w,2,0)\)，并令竖直方向为 \(z\) 轴. 则平面 \(AA_1B_1B\) 为平面 \(x=0\)，且 \(\overrightarrow{BD}=(w,2,0)\). 由直线与平面所成角的定义，
\(\sin30^{\circ}=\dfrac{w}{\sqrt{w^2+4}}\)，解得 \(w=\dfrac2{\sqrt3}\). 此时 \(A_1=(0,2,1)\)，\(B_1=(0,0,1)\)，故 \(F=(0,1,1)\).

\begin{enumerate}
\item 设 \(E=t\overrightarrow{BD}\). 由 \(AE\perp BD\)，有 \((A-E)\cdot\overrightarrow{BD}=0\)，从而 \(t=\dfrac34\)，所以 \(E=\left(\dfrac{3w}{4},\dfrac32,0\right)\). 因此 \(\overrightarrow{AE}=\left(\dfrac{3w}{4},-\dfrac12,0\right)\)，\(\overrightarrow{BF}=(0,1,1)\). 于是
\(\left|\overrightarrow{AE}\cdot\overrightarrow{BF}\right|=\dfrac12\)，\(\left|\overrightarrow{AE}\right|=1\)，\(\left|\overrightarrow{BF}\right|=\sqrt2\). 设所成锐角为 \(\theta\)，则 \(\cos\theta=\dfrac{1/2}{\sqrt2}=\dfrac{\sqrt2}{4}\)，所以 \(\theta=\arccos\dfrac{\sqrt2}{4}\).
\item 平面 \(BDF\) 的一个法向量为
\(\overrightarrow{BD}\times\overrightarrow{BF}=(2,-w,w)\)，平面 \(AA_1B\) 的一个法向量为 \((1,0,0)\). 所以两平面所成锐角 \(\varphi\) 满足 \(\cos\varphi=\dfrac2{\sqrt{4+2w^2}}=\sqrt{\dfrac35}\)，故 \(\varphi=\arccos\sqrt{\dfrac35}\).
\item 由法向量可知平面 \(BDF\) 的方程为 \(2x-wy+wz=0\). 点 \(A=(0,2,0)\) 到该平面的距离为
\(\dfrac{|0-2w+0|}{\sqrt{4+2w^2}}=\dfrac2{\sqrt5}\).
\end{enumerate}
\end{solution}

\begin{problem}
已知数列 \(\{a_{n}\}\) 的首项 \(a_1 = 5\)，前 \(n\) 项和为 \(S_{n}\)，且 \(S_{n+1} = 2S_{n} + n + 5\)（\(n \in \mathbf{N}^{*}\)）.

\begin{enumerate}
\item 证明数列 \(\{a_{n} + 1\}\) 是等比数列；
\item 令 \( f(x) = a_1x + a_2x^2 + \cdots + a_nx^n \)，求函数 \( f(x) \) 在点 \( x = 1 \) 处的导数 \( f'(1) \).
\end{enumerate}
\end{problem}

\begin{solution}
\begin{enumerate}
\item 由已知式，\(a_{n+1}=S_{n+1}-S_n=S_n+n+5\). 当 \(n=1\) 时，\(S_2=2S_1+1+5=16\)，所以 \(a_2=S_2-S_1=11=2a_1+1\). 当 \(n\geqslant2\) 时，\(S_n=2S_{n-1}+n+4\)，从而
\(a_{n+1}+1=S_n+n+6=2S_{n-1}+2n+10=2\left(S_{n-1}+n+5\right)=2(a_n+1)\).
因此对所有 \(n\in\N^*\)，都有 \(a_{n+1}+1=2(a_n+1)\)，且 \(a_1+1=6\)，所以 \(\{a_n+1\}\) 是首项为 \(6\)、公比为 \(2\) 的等比数列.
\item 由上题，\(a_k=6\cdot2^{k-1}-1=3\cdot2^k-1\). 于是
\(f'(1)=\sum_{k=1}^nka_k=3\sum_{k=1}^nk2^k-\sum_{k=1}^nk\). 又 \(\sum_{k=1}^nk2^k=(n-1)2^{n+1}+2\)，\(\sum_{k=1}^nk=\dfrac{n(n+1)}2\)，所以
\(f'(1)=3(n-1)2^{n+1}+6-\dfrac{n(n+1)}2\).
\end{enumerate}
\end{solution}

\begin{problem}
已知动圆过定点 \(\left(\frac{p}{2}, 0\right)\)，且与直线 \(x = -\frac{p}{2}\) 相切，其中 \(p > 0\).

\begin{enumerate}
\item 求动圆圆心 \(C\) 的轨迹的方程；
\item 设 \(A\)，\(B\) 是轨迹 \(C\) 上异于原点 \(O\) 的两个不同点，直线 \(OA\) 和 \(OB\) 的倾斜角分别为 \(\alpha\) 和 \(\beta\)，当 \(\alpha\)，\(\beta\) 变化且 \(\alpha + \beta = \frac\pi4\)，证明直线 \(AB\) 恒过定点，并求出该定点的坐标.
\end{enumerate}
\end{problem}

\begin{solution}
\begin{enumerate}
\item 设动圆圆心为 \(C=(x,y)\)，则圆经过 \(P=\left(\dfrac p2,0\right)\)，且与直线 \(x=-\dfrac p2\) 相切，所以半径的平方分别为
\(\left(x-\dfrac p2\right)^2+y^2\) 和 \(\left(x+\dfrac p2\right)^2\). 令二者相等，得 \(y^2=2px\). 因此圆心轨迹为抛物线 \(y^2=2px\).
\item 设 \(u=\cot\alpha\)，\(v=\cot\beta\). 因为 \(A\)、\(B\) 异于原点，且抛物线 \(y^2=2px\) 与 \(x\) 轴只有原点交点，所以 \(A\)、\(B\) 的纵坐标均不为零. 直线 \(OA\) 的方程可写为 \(x=uy\)，代入抛物线方程得 \(y=2pu\)，从而 \(A=(2pu^2,2pu)\)；同理 \(B=(2pv^2,2pv)\). 直接将这两点代入可得直线 \(AB\) 的方程为
\(x-(u+v)y+2puv=0\). 由 \(\tan(\alpha+\beta)=\tan\dfrac\pi4=1\)，有
\(\dfrac{u+v}{uv-1}=1\)，即 \(u+v=uv-1\). 令 \(s=u+v\)，则 \(uv=s+1\)，直线方程化为
\(x+2p+s(2p-y)=0\). 因此点 \((-2p,2p)\) 对任意允许的 \(u\)，\(v\) 都在直线 \(AB\) 上，故直线 \(AB\) 恒过定点 \((-2p,2p)\).
\end{enumerate}
\end{solution}
